\section{Experimental design}
\label{sec:experimental-design}

\subsection{Reproductions and predictions}

Two kinds of statements are tested, and they differ in what they can show. The first kind restates values that the product committed before this study. The multipliers of the Adverse and Extreme bases, the refusal of the fast-run base, the modelled loss of 0.1708 against 0.85 in two days, and the round trip were written into the committed documents of case C16 on 19 October 2026, and the design document of the product named the same values as the evidence for the proposed method on 21 October 2026, after the documents existed. The note in \texttt{experiments/hypotheses.json} gives 28 October, the date on which the design document's table of methods gained the row for this method; the values themselves were first written into the design document on 21 October. The note is left as it was written, so that its digest in \texttt{experiments/file\_times.json} still matches. These are reproductions, listed as R1 to R5: the engine is called on the committed ladders and must return the committed numbers. They cannot fail on the committed documents and they show that the experiments of this paper start from the same objects.

The second kind is the set of twelve predictions H1 to H12, written by the author of this paper into a file of the experiment folder before the corpus, theory and sensitivity runs. The order is attested only by the modification times that the files had in the working copy where the experiments were run: the predictions file at \tOrdHyp{} and the three result files at \tOrdC{}, \tOrdT{} and \tOrdK{} on 7 November 2026 (UTC$-3$), recorded with the digests of the files in \texttt{experiments/file\_times.json} by \texttt{experiments/record\_order.py}. The predictions, the scripts and the results are in one commit, and a clone does not keep file times, so a reader cannot verify the order from the repository. The file times are those at which each run wrote its result: the corpus run, with its stored wall time, started at \tOrdKStart{}, after the predictions, and the starts of the other two runs were not recorded; a later version should commit the predictions before the runs. They are of three families. H1 to H6 concern the solver on the corpus: that its outcome equals the oracle's on every instance (H1), that its multiplier exceeds the oracle's root by at most the rounding step and satisfies the threshold property (H2), that the round trip is exact with a constant capacity in at least 99\% of the solved instances and has a positive gap in at least 5\% of those with a monetisation schedule (H3), that the case grid overshoots the threshold by more than 10\% in at least 30\% of the solved instances (H4), that Brent's method on the survival period returns a multiplier more than $10^{-6}$ above the threshold in at least 10\% (H5), and that the multiplier $\mathrm{SP}(1)/T$ is within 10\% of the threshold in at least 80\% of the solved instances with a constant capacity (H6). H7 to H9 concern the committed bank: that the derivative of \cref{eq:sensitivity} agrees with repeated solves within 1\% (H7), that the Adverse multiplier moves by less than 25\% across the haircut band and the two monetisation bands of the case design (H8), and that placing the uninsured balance with other legal entities makes the solver refuse (H9). H10 to H12 concern the labelled run: that the fast-run base reaches the amber bar of 0.5 of the observed two-day outflow (H10), that the multiplier of the labelled run is at least twice the multiplier that brings the survival period of that base to 2 days (H11), and that the median or mean calibration rule instead of the maximum does not reach the amber bar (H12).

H10 rests on a bar of the case design, 0.5 for amber, committed when the case was first computed. Two criteria are read more narrowly than they were written, a reading fixed after the runs: the solved instances of H4 are read as those the grid answers (it also holds on all solved instances, \kHFourAll{}\,\%), and the 2 days of H11 as at most 2 days, the target of the solver, at which the survival period is \spAtTwo{} day. The predictions were not tuned after the runs: three of them failed, and \cref{tab:hypotheses} lists the outcomes.

\subsection{Data and provenance}

Two data sources are used. The first is the committed documents of case C16 of the product, a bank whose call-report aggregates are real and whose deposit behaviour is synthetic. The balance sheet is the FDIC call report of Silicon Valley Bank (certificate 24735) at 31 December 2022, with the reports of 31 December 2021 and 30 June 2022 as a time block, each read only up to its date; the haircuts are one minus the Federal Reserve discount window margins in force on 10 March 2023; the run-off rates are those of the supervisory stress test transcribed from its methodological note, on a stated United States business-day calendar; the Table A.1 rows are those of the Basel Committee's report on the 2023 turmoil; and the event amounts are those of the Federal Reserve's review \citep{ecblist2019,bcbs555,fedsvb2023}. The deposit mapping, the monetisation schedule (50\%, 80\% and 100\% of the liquidity value by days 1, 2 and 3), the day calendar and the calibration rule are conventions, named as such in every document, and the raw files are in the product's data root with their hashes recorded. The three bases are LiST Adverse and LiST Extreme (125 relevant days) and the 2023 fast run (5 days). The experiments rebuild each base from the committed ladder, with the scenario, the effective capacity and the engine calls identical to those of the case pipeline, and check the rebuilt multipliers and curves against the committed document before any sensitivity is computed.

The second source is a corpus of synthetic banks with known structure, generated by committed code from seeds. A bank has two to six deposit segments, a financial segment of rate 1 with probability 0.15, constant daily rates drawn log-uniformly between 0.0005 and 0.08 (with a two-phase path that starts faster for the first two to five days with probability 0.3), a capacity of 3\% to 50\% of the deposits released by an immediate, a three-day or a ten-day monetisation schedule, a day-1 inflow with probability 0.3 and a constant daily outflow outside the segments with probability 0.5, and a horizon of 5, 30, 60 or 125 days. The rates cover the supervisory values of the committed bank, which range from 0.001 to 0.011 per day for the non-financial classes. The corpus is constructed so that the three refusal reasons occur, and the reader can see the mix in \cref{tab:design}.

\subsection{Methods compared}

The proposed solver is the engine's \texttt{solve\_multiplier}. The comparators are the ways a validator who has the forward tools of the supervisory exercise would invert them, and each is run by the experiment code on the same instances. The first is a scan of the multiplier grid of the case design (1, 1.25, 1.5, 2, 3, 5, 8, 13, 21, 34, 55, 89, 144, truncated at the cap) with the engine's forward function \texttt{survival\_vs\_multiplier}, taking the first multiplier at which the survival period is at most the target; a scan of a finer grid (step 0.1 up to 20 and the case grid above it) shows what the extra evaluations buy. The third is Brent's method on $\mathrm{SP}(m)-T$, a survival that never turns negative counted as $H+1$, the literal reading of the multiplier at which the survival period equals the target. The fourth is a bisection of the predicate that the target is reached on a fixed bracket $[1,100]$ that returns the right end when the bracket fails and the left end when it holds at 1. Three further variants are ablations of the proposed solver: without the cap, without the upward rounding and neighbour check, and handed the capacity schedule alone, without the inflows and the flows outside the segments.

The oracle shares the model of \cref{eq:outflow,eq:nlp-sp} and none of the solver code. It computes the position by the closed form of the stock, finds the root of $w_T$ by two hundred bisections to a relative width of $10^{-13}$, and applies the three refusal conditions in the order of the rule.

\subsection{Metrics}

For an instance with oracle outcome $o\in\{\text{solved},\text{refused}\}$ and a method that either answers with a number $\hat m$ or does not, the experiments report the share of refused instances that are answered, $\Pr(\text{answer}\mid\text{refused})$, the share of solved instances without an answer, $\Pr(\text{no answer}\mid\text{solved})$, the relative error and the threshold property,
\begin{equation}
e=\frac{\hat m-m_0}{m_0},\qquad \tau=\mathbf 1\bigl[\mathrm{SP}(\hat m)\le T<\mathrm{SP}(\hat m-\delta)\bigr],
\label{eq:metrics}
\end{equation}
with $\mathrm{SP}$ evaluated by the independent forward computation and not by the method's own. A fifth metric is the share of all instances answered with a multiplier above the cap, which would drive a rate above 1. The cost is the number of evaluations of the position per solved instance and the measured median time of a call. The error is signed in the stored results, so that a method that systematically overshoots, as a grid scan does, is told apart from one that scatters around the threshold, and the threshold property is evaluated on the independent forward computation so that a method cannot pass by reporting its own survival period. For the committed bank the metrics are the multiplier, the survival periods at and just below it, the change under each convention, and for the labelled run the ratio of the modelled to the observed share.

\subsection{Protocol and statistics}

Five independent corpora, seeds 20261107 to 20261111, of 1,000 banks each are solved at the targets of 2, 5, 10 and 30 days. Rates are reported on the pooled counts with 95\% Wilson intervals \citep{brown2001} and with the standard deviation of the rate over the five corpora; quantiles of the relative error are reported with bootstrap percentile intervals of 1,000 resamples \citep{efron1979}; the spread and the intervals are given in \cref{tab:spread}. The theory checks use their own seed, 20261120. The design follows the guidance for simulation studies on stating the aims, the data-generating mechanism, the methods, the estimand and the performance measures \citep{morris2019}. The five corpora are independent draws on disjoint seeds, so the spread of a rate over them measures the variation of the corpus and not of the solver, which is deterministic; the Wilson interval on the pooled counts treats the instances as independent, which they are not exactly, since the four targets of a bank share its structure, and the spread over corpora is the check on that. The experiments are run once per seed and the stored results are the records, because the solver and the oracle are deterministic and a repetition would reproduce them. The committed bank is a single deterministic case, so its uncertainty is the range over the conventions of the case design and the two denominators of the observed outflow, not a sampling interval.
